% ============================================================================
% Fichier  : partie2/P02_C05_politique.tex
% Titre    : La Politique de Sécurité du Système d'Information
% Auteur   : MINKA MI NGUIDJOÏ Thierry Emmanuel
% ============================================================================
% Sources : ISO/IEC 27001:2022 (clause 5.2) ; ISO/IEC 27002:2022. Fictif : Eséka.
% ============================================================================

\chapter{La Politique de Sécurité du Système d'Information}
\label{chap:politique}
\chaptermark{La politique de sécurité}

\epigraph{%
    « Une politique qui ne s'applique pas est pire que pas de politique~:\\
    elle donne l'illusion d'être protégé. »%
}{Principe du cours}

\niveauChapitre{Fondamental}{Le Chapitre~\ref{chap:gouvernance}.}

\begin{boxobjectifs}
\begin{itemize}
    \item Définir la politique de sécurité (\sigle{PSSI}) et ses objectifs.
    \item Décrire la structure type d'une politique de sécurité.
    \item Distinguer politique, procédure et charte.
    \item Comprendre le cycle de vie d'une politique (validation, diffusion,
          révision).
    \item Relier la politique à l'exigence de la clause~5.2 de la norme.
\end{itemize}
\end{boxobjectifs}

\bigskip

% ============================================================================
\section*{Question directrice : comment transformer des intentions en règles opposables ?}
% ============================================================================

Le maire veut « plus de sécurité ». Cette intention ne protège rien tant
qu'elle n'est pas écrite, validée, diffusée et appliquée. La politique de
sécurité est le document qui opère cette transformation.

\medskip
\textbf{Plan du chapitre.}
\begin{enumerate}
    \item Qu'est-ce qu'une politique (section~\ref{sec:c05-definir}).
    \item Objectifs et portée (section~\ref{sec:c05-objectifs}).
    \item La structure type (section~\ref{sec:c05-structure}).
    \item Politique, procédure, charte (section~\ref{sec:c05-niveaux}).
    \item Le cycle de vie du document (section~\ref{sec:c05-cycle}).
\end{enumerate}

\separateur

% ============================================================================
\section{Qu'est-ce qu'une Politique}
\label{sec:c05-definir}
% ============================================================================

\begin{boxdef}
\termecle{Politique de sécurité du système d'information} (\sigle{PSSI})~:
document formel, approuvé par la direction, qui fixe les objectifs de sécurité,
les règles à respecter et les responsabilités. Elle est le cadre de référence
de toutes les décisions de sécurité.
\end{boxdef}

La \sigle{PSSI} est au \sigle{SMSI} ce que la constitution est à un État~: le
texte fondateur dont découlent les règles particulières. Elle doit être courte,
claire, et tenue par la direction, sans quoi elle reste lettre morte.

% ============================================================================
\section{Objectifs et Portée}
\label{sec:c05-objectifs}
% ============================================================================

\begin{tableaucours}{L{4.0cm}L{8.2cm}}{%
    \tableauHeader{Objectif} &
    \tableauHeader{Ce qu'il apporte}
}
Établir un cadre & Une référence commune pour arbitrer \\
Sensibiliser & Des règles connues de tous les agents \\
Définir les responsabilités & Chacun sait ce qu'il doit faire \\
Assurer la conformité & Le respect des obligations légales \\
Faciliter la décision & Un critère pour trancher les cas particuliers \\
\end{tableaucours}
\vspace{-0.8em}
\begin{center}{\small\itshape Tableau~5.1. Les objectifs d'une politique de sécurité}\end{center}

La \termecle{portée} délimite ce que couvre la politique~: quels services,
quels sites, quelles informations. À Eséka, la politique couvre tous les
services de la mairie et l'ensemble du patrimoine informationnel recensé.

% ============================================================================
\section{La Structure Type}
\label{sec:c05-structure}
% ============================================================================

\begin{tableaucours}{L{4.2cm}L{8.0cm}}{%
    \tableauHeader{Section} &
    \tableauHeader{Contenu}
}
Introduction et contexte & Objectifs, portée, références réglementaires \\
Gouvernance & Organisation, rôles, comité de sécurité \\
Classification de l'information & Niveaux et règles de manipulation \\
Contrôle d'accès & Création des comptes, authentification, privilèges \\
Sécurité des opérations & Sauvegarde, journalisation, gestion des incidents \\
Conformité & Audits, revue, sanctions \\
\end{tableaucours}
\vspace{-0.8em}
\begin{center}{\small\itshape Tableau~5.2. Structure type d'une politique de sécurité}\end{center}

\begin{boxmetier}{RSSI}
\textbf{Écrire pour être lu.} La RSSI d'Eséka proscrit le jargon~: une
politique que les agents ne comprennent pas n'est pas appliquée. Chaque règle
dit qui, quoi, quand, et la conséquence du non-respect.
\end{boxmetier}

% ============================================================================
\section{Politique, Procédure, Charte}
\label{sec:c05-niveaux}
% ============================================================================

\begin{boxdef}
Trois niveaux se complètent~: la \termecle{politique} fixe les principes (le
« quoi »)~; la \termecle{procédure} décrit les étapes d'exécution (le
« comment »)~; la \termecle{charte} engage l'utilisateur sur ses droits et
devoirs.
\end{boxdef}

La politique dit « les accès sont nominatifs »~; la procédure décrit comment
créer et supprimer un compte~; la charte informatique, signée par l'agent,
l'engage à ne pas partager son mot de passe (Chapitre~\ref{chap:conformite}).

% ============================================================================
\section{Le Cycle de Vie du Document}
\label{sec:c05-cycle}
% ============================================================================

Une politique vit~: elle est \termecle{rédigée}, \termecle{validée} par la
direction, \termecle{diffusée} à tous, \termecle{appliquée}, puis
\termecle{révisée} périodiquement. Sans validation de la direction, elle n'a
pas d'autorité~; sans révision, elle se périme.

\begin{boxprocedure}
\textbf{Constatation.} La mairie a rédigé une politique en 2022, jamais signée
par le maire ni communiquée aux agents.

\medskip
\textbf{Analyse} (confiance élevée). Un document non validé et non diffusé
n'est pas une politique~: c'est un brouillon. L'opération \textsc{Nyong}
reprendra le cycle complet, de la validation à la diffusion contre signature.
\end{boxprocedure}

\begin{boxjuris}
La clause~5.2 de \norme{ISO/IEC 27001} exige une politique de sécurité
établie, approuvée par la direction, communiquée et disponible. L'absence de
politique formelle est un écart majeur en audit de certification
(Chapitre~\ref{chap:audit}).
\end{boxjuris}

\separateur

\begin{boxresume}
\textbf{Ce qu'il faut retenir de ce chapitre~:}
\begin{itemize}
    \item La \textbf{PSSI} est le texte fondateur du \sigle{SMSI}, approuvé par
          la direction.
    \item Elle poursuit cinq \textbf{objectifs} et définit une \textbf{portée}.
    \item Sa \textbf{structure type} couvre gouvernance, accès, opérations,
          conformité.
    \item \textbf{Politique, procédure, charte} se distinguent par leur niveau.
    \item Une politique a un \textbf{cycle de vie}~: validée, diffusée,
          révisée~; la clause~5.2 l'exige.
\end{itemize}
\end{boxresume}

\begin{boxinfo}
\textbf{Points de contrôle de l'élève.} Je sais~: définir la \sigle{PSSI}~;
citer ses objectifs~; décrire sa structure~; distinguer politique, procédure et
charte~; énoncer son cycle de vie.
\end{boxinfo}

\bigskip

\begin{boxexo}[title={\ding{45}\quad Exercice~5.1. Définition \niveaubase}]
Qu'est-ce qu'une politique de sécurité, et qui doit l'approuver~?
\end{boxexo}

\begin{boxexo}[title={\ding{45}\quad Exercice~5.2. Objectifs \niveaubase}]
Citez trois objectifs d'une \sigle{PSSI} et illustrez-en un pour la mairie.
\end{boxexo}

\begin{boxexo}[title={\ding{45}\quad Exercice~5.3. Trois niveaux \niveaubase}]
Classez~: « les accès sont nominatifs »~; « voici comment créer un compte »~;
« je m'engage à ne pas partager mon mot de passe ». Politique, procédure ou
charte~?
\end{boxexo}

\begin{boxexo}[title={\ding{45}\quad Exercice~5.4. Portée \niveaumoyen}]
Rédigez la phrase de portée de la politique de la mairie d'Eséka.
\end{boxexo}

\begin{boxexo}[title={\ding{45}\quad Exercice~5.5. Brouillon ou politique \niveaumoyen}]
Pourquoi le document de 2022 n'est-il pas une politique~? Que manque-t-il~?
\end{boxexo}

\begin{boxexo}[title={\ding{45}\quad Exercice~5.6. Une règle complète \niveaumoyen}]
Rédigez une règle de politique sur l'usage des clés \texttt{USB}, avec la
conséquence du non-respect.
\end{boxexo}

\begin{boxexo}[title={\ding{45}\quad Exercice~5.7. Cycle de vie \niveaumoyen}]
Énoncez les cinq étapes du cycle de vie d'une politique et le risque d'en
oublier une.
\end{boxexo}

\begin{boxexo}[title={\ding{45}\quad Exercice~5.8. Écart d'audit \niveaumoyen}]
Pourquoi l'absence de politique validée est-elle un écart majeur en audit~?
\end{boxexo}

\begin{boxexo}[title={\ding{45}\quad Exercice~5.9. Table des matières \niveauexpert}]
Dressez la table des matières complète de la politique de sécurité de la mairie
d'Eséka.
\end{boxexo}

\begin{boxexo}[title={\ding{45}\quad Exercice~5.10. Section rédigée \niveauexpert}]
Rédigez la section « Contrôle d'accès » de la politique de la mairie (une
page).
\end{boxexo}

\bigskip

\begin{boxinfo}
\textbf{Pour aller plus loin.}
\begin{itemize}
    \item ISO/IEC 27001, clause~5.2~\cite{iso27001}~; ISO/IEC 27002~\cite{iso27002}.
\end{itemize}
\end{boxinfo}

\begin{boxinfo}
\textbf{Transition.} La politique s'adosse à des obligations légales et à des
engagements individuels~: le Chapitre~\ref{chap:conformite} traite chartes,
conformité et données personnelles.
\end{boxinfo}
